%% bare_jrnl.tex
%% V1.4b
%% 2015/08/26
%% by Michael Shell
%% see http://www.michaelshell.org/
%% for current contact information.
%%
%% This is a skeleton file demonstrating the use of IEEEtran.cls
%% (requires IEEEtran.cls version 1.8b or later) with an IEEE
%% journal paper.
%%
%% Support sites:
%% http://www.michaelshell.org/tex/ieeetran/
%% http://www.ctan.org/pkg/ieeetran
%% and
%% http://www.ieee.org/

%%*************************************************************************
%% Legal Notice:
%% This code is offered as-is without any warranty either expressed or
%% implied; without even the implied warranty of MERCHANTABILITY or
%% FITNESS FOR A PARTICULAR PURPOSE! 
%% User assumes all risk.
%% In no event shall the IEEE or any contributor to this code be liable for
%% any damages or losses, including, but not limited to, incidental,
%% consequential, or any other damages, resulting from the use or misuse
%% of any information contained here.
%%
%% All comments are the opinions of their respective authors and are not
%% necessarily endorsed by the IEEE.
%%
%% This work is distributed under the LaTeX Project Public License (LPPL)
%% ( http://www.latex-project.org/ ) version 1.3, and may be freely used,
%% distributed and modified. A copy of the LPPL, version 1.3, is included
%% in the base LaTeX documentation of all distributions of LaTeX released
%% 2003/12/01 or later.
%% Retain all contribution notices and credits.
%% ** Modified files should be clearly indicated as such, including  **
%% ** renaming them and changing author support contact information. **
%%*************************************************************************


% *** Authors should verify (and, if needed, correct) their LaTeX system  ***
% *** with the testflow diagnostic prior to trusting their LaTeX platform ***
% *** with production work. The IEEE's font choices and paper sizes can   ***
% *** trigger bugs that do not appear when using other class files.       ***                          ***
% The testflow support page is at:
% http://www.michaelshell.org/tex/testflow/



\documentclass[journal]{IEEEtran}
%
% If IEEEtran.cls has not been installed into the LaTeX system files,
% manually specify the path to it like:
% \documentclass[journal]{../sty/IEEEtran}





% Some very useful LaTeX packages include:
% (uncomment the ones you want to load)


% *** MISC UTILITY PACKAGES ***
%
%\usepackage{ifpdf}
% Heiko Oberdiek's ifpdf.sty is very useful if you need conditional
% compilation based on whether the output is pdf or dvi.
% usage:
% \ifpdf
%   % pdf code
% \else
%   % dvi code
% \fi
% The latest version of ifpdf.sty can be obtained from:
% http://www.ctan.org/pkg/ifpdf
% Also, note that IEEEtran.cls V1.7 and later provides a builtin
% \ifCLASSINFOpdf conditional that works the same way.
% When switching from latex to pdflatex and vice-versa, the compiler may
% have to be run twice to clear warning/error messages.






% *** CITATION PACKAGES ***
%
\usepackage{cite}
% cite.sty was written by Donald Arseneau
% V1.6 and later of IEEEtran pre-defines the format of the cite.sty package
% \cite{} output to follow that of the IEEE. Loading the cite package will
% result in citation numbers being automatically sorted and properly
% "compressed/ranged". e.g., [1], [9], [2], [7], [5], [6] without using
% cite.sty will become [1], [2], [5]--[7], [9] using cite.sty. cite.sty's
% \cite will automatically add leading space, if needed. Use cite.sty's
% noadjust option (cite.sty V3.8 and later) if you want to turn this off
% such as if a citation ever needs to be enclosed in parenthesis.
% cite.sty is already installed on most LaTeX systems. Be sure and use
% version 5.0 (2009-03-20) and later if using hyperref.sty.
% The latest version can be obtained at:
% http://www.ctan.org/pkg/cite
% The documentation is contained in the cite.sty file itself.

\usepackage{fnpct}




% *** GRAPHICS RELATED PACKAGES ***
%
\ifCLASSINFOpdf
  % \usepackage[pdftex]{graphicx}
  % declare the path(s) where your graphic files are
  % \graphicspath{{../pdf/}{../jpeg/}}
  % and their extensions so you won't have to specify these with
  % every instance of \includegraphics
  % \DeclareGraphicsExtensions{.pdf,.jpeg,.png}
\else
  % or other class option (dvipsone, dvipdf, if not using dvips). graphicx
  % will default to the driver specified in the system graphics.cfg if no
  % driver is specified.
  % \usepackage[dvips]{graphicx}
  % declare the path(s) where your graphic files are
  % \graphicspath{{../eps/}}
  % and their extensions so you won't have to specify these with
  % every instance of \includegraphics
  % \DeclareGraphicsExtensions{.pdf}
\fi
% graphicx was written by David Carlisle and Sebastian Rahtz. It is
% required if you want graphics, photos, etc. graphicx.sty is already
% installed on most LaTeX systems. The latest version and documentation
% can be obtained at: 
% http://www.ctan.org/pkg/graphicx
% Another good source of documentation is "Using Imported Graphics in
% LaTeX2e" by Keith Reckdahl which can be found at:
% http://www.ctan.org/pkg/epslatex
%
% latex, and pdflatex in dvi mode, support graphics in encapsulated
% postscript (.pdf) format. pdflatex in pdf mode supports graphics
% in .pdf, .jpeg, .png and .mps (metapost) formats. Users should ensure
% that all non-photo figures use a vector format (.pdf, .pdf, .mps) and
% not a bitmapped formats (.jpeg, .png). The IEEE frowns on bitmapped formats
% which can result in "jaggedy"/blurry rendering of lines and letters as
% well as large increases in file sizes.
%
% You can find documentation about the pdfTeX application at:
% http://www.tug.org/applications/pdftex





% *** MATH PACKAGES ***
%
\usepackage{amsmath}
\allowdisplaybreaks
\usepackage{amsfonts}
\usepackage{mathtools}
\usepackage{amsthm}
\theoremstyle{definition}
\newtheorem{definition}{Definition}
\newtheorem{prop}{Proposition}
\newtheorem{theorem}{Theorem}

\theoremstyle{remark}
\newtheorem{remark}{Remark}
\usepackage[ruled,norelsize,linesnumbered]{algorithm2e}

\makeatletter
\newcommand{\removelatexerror}{\let\@latex@error\@gobble}
\makeatother
\usepackage{color}
\usepackage[normalem]{ulem}
\usepackage{tabularray}
\usepackage{booktabs}
\usepackage[flushleft]{threeparttable}
\usepackage[caption=false,font=footnotesize]{subfig}
\usepackage[utf8]{inputenc}
\usepackage[T1]{fontenc}
\usepackage{tikz,pgfplots}
\usetikzlibrary{positioning}

\hyphenation{op-tical net-works semi-conduc-tor}


\begin{document}

\title{Computationally Enhanced Approach for Chance-Constrained OPF Considering \\Voltage Stability}


\author{Yuanxi~Wu,~\IEEEmembership{Graduate~Student~Member,~IEEE,}
        Zhi~Wu,~\IEEEmembership{Senior~Member,~IEEE,}
        Yijun~Xu,~\IEEEmembership{Senior~Member,~IEEE,}
        Huan~Long,\IEEEmembership{Member,~IEEE,}
        Wei~Gu,~\IEEEmembership{Senior~Member,~IEEE,}
        Shu~Zheng,
        and~Jingtao~Zhao% <-this % stops a space
        \thanks{This work has been submitted to the IEEE for possible publication. Copyright may be transferred without notice, after which this version may no longer be accessible.}
        \thanks{The study was supported by the science and technology project of the State Grid Corporation of China, Grant 5108-202218280A-2-366-XG.}
\thanks{Yuanxi Wu, Zhi Wu, Yijun Xu, Huan Long, Wei Gu, and Shu Zheng are with the School
of Electrical Engineering, Southeast University, Nanjing 210096, China (e-mail: yuanxi\underline{~~}wu@seu.edu.cn; zwu@seu.edu.cn; yijunxu@seu.edu.cn; hlong@seu.edu.cn; wgu@seu.edu.cn; zhengshu@sgepri.sgcc.com.cn).}% <-this % stops a space
\thanks{Jingtao Zhao is with State Grid Electric Power Research Institute, Nanjing 211102, China (e-mail: zhaojingtao@sgepri.sgcc.com.cn).}% <-this % stops a space
}

 \markboth{SUBMITTED FOR REVIEW}%
 {Shell \MakeLowercase{\textit{et al.}}: Bare Demo of IEEEtran.cls for IEEE Journals}

\maketitle

\begin{abstract}
The effective management of stochastic characteristics of renewable power generations is vital for ensuring the stable and secure operation of power systems. This paper addresses the formidable task of optimizing the chance-constrained voltage-stability-constrained optimal power flow (CC-VSC-OPF) problem, which is hindered by the implicit voltage stability index and intractable chance constraints.
Leveraging a neural network (NN)-based surrogate model, the stability constraint is explicitly formulated and seamlessly integrated into the model. To perform uncertainty propagation without relying on presumptions or complicated transformations, an advanced data-driven method known as adaptive polynomial chaos expansion (APCE) is developed. To extend the scalability of the proposed algorithm, a partial least squares (PLS)-NN framework is designed, which enables the establishment of a parsimonious surrogate model and efficient computation of large-scale Hessian matrices. In addition, a dimensionally decomposed APCE (DD-APCE) is proposed to alleviate the "curse of dimensionality" by restricting the interaction order among random variables. Finally, the above techniques are merged into an iterative scheme to update the operation point.
Simulation results reveal the cost-effective performances of the proposed method in several test systems. 
\end{abstract}

\begin{IEEEkeywords}
Optimal power flow, chance constraint, PLS, neural network, dimensionally decomposed APCE, data-driven.
\end{IEEEkeywords}

\IEEEpeerreviewmaketitle



\section{Introduction}\label{introduction}
\IEEEPARstart{A}{s} the world shifts up a gear in the development of renewable energy resources, uncertainties in power systems inevitably arise since renewable generations can hardly be forecasted with high accuracy \cite{kirby2017analyzing}.
Deterministic AC optimal power flow (OPF), although widely adopted in the traditional power industry to achieve reliable system operations, is not intended for such circumstances. Failing to incorporate stochastic features would result in significant degradation of system performance \cite{wang2017efficient}, which urges the necessity of moving beyond this classical setting.  

Different approaches exist to immunize power systems against uncertain injections. One intensively studied type is robust optimization \cite{qiu2020historical}, which enforces the power system to remain on the safe side for all possible uncertain-but-bounded perturbations. However, this worst-case-oriented philosophy is inclined to deliver an over-conservative solution, thus rendering itself less preferable. Much recent research effort has been directed to chance-constrained AC-OPF (CC-AC-OPF), which requires the uncertainty-affected constraints to be satisfied with a predefined probability. While capable of promoting operational efficiency, CC-AC-OPF is, more often than not, a heavily problematic task \cite{pena2020solving}. The crux of its solution lies in the \textit{ efficient evaluation of chance constraints}, which constitutes the first research concern of this paper. 

In light of this challenge, the traditional Monte-Carlo (MC) simulation offers a straightforward and conceptually clear-cut solution \cite{zhang2010probabilistic}. However, this method is extremely time-consuming, even for moderately sized problems. Du \textit{et al.} \cite{du2021chance} adopt the point estimation method to avoid massive sampling. Notwithstanding the improved efficiency, the accuracy of this method is not guaranteed, and its application to non-Gaussian random variables with nonlinear correlations is quite limited. 

Alternatively, the analytical method offers a viable approach for transforming chance constraints into tractable formulations. Early research relies on strong assumptions regarding characteristics of renewable generations, such as Gaussianity and independence. Instead, Yang \textit{et al.} \cite{yang2020analytical} embrace the adoption of the Gaussian mixture model to enable general characterizations of non-Gaussian uncertainties. Another widely adopted technique is the distributionally robust model, which utilizes certain statics of uncertainties \cite{roald2018chance,li2019distribution,arab2022distributionally,arrigo2022embed}. It is notable, however, that most of the studies cited above resort to approximations of AC power flow (e.g., linearized or relaxed power flow).

Recently, meta-modeling techniques have received significant attention due to their ability to inexpensively deal with uncertainty in AC power systems. This paper specifically focuses on variants of polynomial chaos expansion (PCE). Despite previous successful attempts to apply PCE to CC-OPF \cite{li2018compressive,muhlpfordt2019chance,metivier2020efficient,xu2021iterative}, a common trait of these works is that they have simplified the underlying dependence structure of random variables to varying degrees.
While recent studies in probabilistic power flow have made progress via the copula technique and Rosenblatt transformation
 \cite{ye2022generalized,ly2022scalable}, selecting an appropriate copula for modeling the joint probability distribution can prove difficult. Furthermore, the highly nonlinear Rosenblatt transformation risks compromising the performance of PCE in practical applications \cite{lee2022reliability}. A promising avenue for improving PCE is the direct construction of multivariate orthogonal polynomials with respect to the original probabilistic space. To the best of our knowledge, this approach is yet to receive sufficient exploration.

 In addition, voltage stability is crucial to ensure the robustness of OPF solutions against system instability. Traditionally, voltage-stability-constrained OPF (VSC-OPF) has been the preferred method to accomplish this goal \cite{jia2023voltage,cui2018new,avalos2008practical}. VSC-OPF enforces a predetermined threshold on the voltage stability index (VSI), typically represented by the minimum singular value of the power flow Jacobian matrix. In the context of high renewable penetration, the system state becomes more volatile, leading to an increased risk of static voltage instability.
 However, incorporating voltage stability constraints becomes more intricate when VSI is subjected to uncertainties. Selecting an appropriate operation point to ensure the probability distribution function of VSI meets the secure operation requirement is a challenging problem.
 Currently, there is very limited research on stability issues in CC-OPF problems, with only a few studies considering small-signal stability \cite{wang2023stability} or static voltage stability margin \cite{yang2022optimization}. 
 Hence, further investigation is necessary to advance the field of \textit{CC-VSC-OPF}, which constitutes the second research concern of this paper. 

In response to these pressing concerns, this study applies a novel adaptive polynomial chaos expansion (APCE) \cite{lee2020practical} to the CC-VSC-OPF problem, where a neural network (NN) regression model is adopted as a surrogate for the voltage stability constraint. The proposed APCE method facilitates stable computation of orthonormal polynomial basis even under arbitrarily distributed and complexly correlated uncertainties. These basis functions are ultimately integrated into an iterative framework to obtain the optimal operation point. Partial least squares (PLS) and a dimensionwise decomposition technique are further incorporated to improve the scalability of the proposed approach. The contributions of this study are summarized as follows:
\begin{itemize}
\item A novel CC-VSC-OPF model is designed with the consideration of voltage stability constraint and uncertain renewable generations, thereby maintaining stable and secure operations of power systems. 
\item This study merges the NN-based surrogate model and APCE into an iterative scheme to deal with the implicit VSI and chance constraints, which transform the original intractable model into standard nonlinear problems.
\item The efficiency and scalability of the proposed algorithm are extended by introducing the PLS-NN framework, which provides a practical workaround to numerically compute the large-scale Hessian matrix with negligible computation burden. A dimensionally decomposed APCE (DD-APCE) is further devised to enable the implementation of the proposed algorithm under high-dimensional uncertainties. 
\end{itemize}
The remainder of this paper is organized as follows: Section \ref{preliminary} introduces the formulation of CC-VSC-OPF. Section \ref{approach} presents the data-driven APCE and its dimensionally decomposed version. Section \ref{scheme} outlines the PLS-NN framework, along with the whole solution procedure for CC-VSC-OPF. Numerical experiments are conducted in Section \ref{simulation}, and conclusions follow in Section \ref{conclusion}.
\section{Preliminaries}\label{preliminary}
This section first outlines the formulation of CC-VSC-OPF by imposing a lower bound on the VSI. Then the NN-based surrogate model is constructed to incorporate the voltage stability constraint into the optimization process explicitly.  
\vspace{5pt}
\subsection{Formulation of CC-VSC-OPF}
Consider an electrical network constituted by a set \(\mathcal{N}\) of \(N=\left|\mathcal{N}\right|\) buses and a set \(\mathcal{L}\) of \(L=\left|\mathcal{L}\right|\) lines. Subscripts $_{PQ}$, $_{PV}$, and $_{ref}$ are added to classify three types of buses, namely PQ, PV, and reference buses. Let $v_i$ and $\theta_i$ denote the voltage magnitude and the voltage angle at each bus. To avoid messy formulas, it is assumed that the power injection at each bus $i\in\mathcal{N}$ can be decomposed into one controllable generation $p_{g,i},q_{g,i}$, one fixed demand $p_{d,i},q_{d,i}$, and one uncertain renewable generation $\xi_{i}$\footnote{The renewable energy stations are assumed to be operated at unity power factor here. However, other types of control can be incorporated as well.}. Consider a random vector $\boldsymbol{\xi}\coloneqq[\xi_1,\xi_2,\dots,\xi_N]^{\mathsf{T}}$ that comprises uncertain renewable generations at each bus. The propagation of these fluctuations through the power system renders the controllable generation dispatch, voltage magnitudes as well as voltage angles functions of $\boldsymbol{\xi}$. In this paper, we assume the active power mismatch is balanced through Automatic Generation Control (AGC). More details can be found in \cite{roald2018chance}. Thus, the net active power injection at bus \textit{i} can be expressed as follows:
\begin{equation}\label{eq:1}
    p_i(\boldsymbol{\xi})=p_{g,i}(\boldsymbol{\xi})+\xi_{i}-p_{d,i},\quad \forall i \in \mathcal{N},
\end{equation}
Similarly, the net reactive power injection at bus \textit{i} is given by:
\begin{equation}\label{eq:2}
    q_i(\boldsymbol{\xi})=q_{g,i}(\boldsymbol{\xi})-q_{d,i},\quad \forall i \in \mathcal{N},
\end{equation}

The electrical network is governed by the AC nodal power balance equations, which are expressed as functions of bus injections and bus voltages in polar coordinates: 
\begin{subequations}\label{eq:4}
    \begin{alignat}{2}        
    p_i(\boldsymbol{\xi})=&v_i(\boldsymbol{\xi})\sum\limits_{j=1}^{N}v_j(\boldsymbol{\xi})\Bigl(G_{ij}\cos(\theta_{i}(\boldsymbol{\xi})-\theta_{j}(\boldsymbol{\xi})) \notag\\
    &+B_{ij}\sin(\theta_{i}(\boldsymbol{\xi})-\theta_{j}(\boldsymbol{\xi}))\Bigr),\quad \forall i \in\mathcal{N}, \\
    q_i(\boldsymbol{\xi})=&v_i(\boldsymbol{\xi})\sum\limits_{j=1}^{N}v_j(\boldsymbol{\xi})\Bigl(G_{ij}\sin(\theta_{i}(\boldsymbol{\xi})-\theta_{j}(\boldsymbol{\xi})) \notag\\
    &-B_{ij}\cos(\theta_{i}(\boldsymbol{\xi})-\theta_{j}(\boldsymbol{\xi}))\Bigr),\quad \forall i \in\mathcal{N},
    \end{alignat}
\end{subequations}
where $G_{ij}$ and $B_{ij}$ represent the real and imaginary components of the bus admittance matrix.

At this point, we have all the ingredients to state the CC-VSC-OPF problem.
\begin{subequations}\label{eq:5}
\allowdisplaybreaks
    \begin{alignat}{3}
        &\!\min &\quad&\sum\limits_{i\in\mathcal{N}_{PV,ref}}\mathbb{E}\left[C_i(p_{g,i})\right]&\quad &  \label{eq:5a}\\
        &\rm{s.t.}  &\quad &\boldsymbol{g}(\boldsymbol{p}(\boldsymbol{\xi}),\boldsymbol{q}(\boldsymbol{\xi}),\boldsymbol{v}(\boldsymbol{\xi}),\boldsymbol{\theta}(\boldsymbol{\xi}))=0,&\quad& \label{eq:5b}\\
        & &\quad&v_i^{\min}\leq\ v_i\leq v_i^{\max}, &\quad&\forall i \in \mathcal{N}_{PV}, \label{eq:5c}\\
        & &\quad&\mathbb{P}(p_{g,i}(\boldsymbol{\xi})\geq p_{g,i}^{\min})\geq 1-\epsilon_P,&\quad &\forall i \in \mathcal{N}_{PV,ref}, \label{eq:5d}\\
        & &\quad&\mathbb{P}(p_{g,i}(\boldsymbol{\xi})\leq p_{g,i}^{\max})\geq 1-\epsilon_P,&\quad &\forall i \in \mathcal{N}_{PV,ref}, \label{eq:5e}\\
        & &\quad &\mathbb{P}(q_{g,i}(\boldsymbol{\xi})\geq q_{g,i}^{\min})\geq 1-\epsilon_Q,&\quad &\forall i \in \mathcal{N}_{PV,ref},\label{eq:5f}\\
        & &\quad &\mathbb{P}(q_{g,i}(\boldsymbol{\xi})\leq q_{g,i}^{\max})\geq 1-\epsilon_Q,&\quad &\forall i \in \mathcal{N}_{PV,ref},\label{eq:5g}\\
        &  &\quad &\mathbb{P}(v_{i}(\boldsymbol{\xi})\geq v_{i}^{\min})\geq 1-\epsilon_V,&\quad &\forall i \in \mathcal{N}_{PQ},\label{eq:5h}\\
        &  &\quad &\mathbb{P}(v_{i}(\boldsymbol{\xi})\leq v_{i}^{\max})\geq 1-\epsilon_V,&\quad &\forall i \in \mathcal{N}_{PQ},\label{eq:5i}\\
        &  &\quad &\mathbb{P}(i_{ij}(\boldsymbol{\xi})\leq i_{ij}^{\max})\geq 1-\epsilon_I, &\quad &\forall ij\in\mathcal{L}, \label{eq:5j}\\
        &  &\quad &\mathbb{P}(\sigma(\boldsymbol{\xi})\geq \sigma^{\min})\geq 1-\epsilon_{\sigma}. &\quad & \label{eq:5k}
    \end{alignat}
\end{subequations}

The objective \eqref{eq:5a} minimizes the total expected cost of active power generation. \eqref{eq:5b} represents the compact form of nodal power balance equations. \eqref{eq:5c} place constraints on the voltage setpoints of PV buses. \eqref{eq:5d}-\eqref{eq:5k} impose chance constraints on other variables subject to uncertainties 
with a desired level of confidence (or reliability)
$1-\epsilon_{P/Q/V/I/\sigma}$ for some small quantities $\epsilon_{P/Q/V/I/\sigma}$ (e.g., 1\% or 5\%). 
\begin{remark}
    Note that \eqref{eq:5k} guarantees a minimum distance from the voltage instability point, where the VSI is denoted by $\sigma(\boldsymbol{\xi})$. It is acknowledged that voltage stability is primarily concerned with the singularity of the power flow Jacobian $\boldsymbol{J}$. Typically, the VSI $\sigma$ is defined as the minimum singular value (MSV) of $\boldsymbol{J}$ since this value tends to zero as the system approaches the voltage collapse point. Despite intuitive, the main drawback of this representation is that the MSV is an implicit function of optimization variables, which renders the constraint \eqref{eq:5k} challenging to be directly embedded into the OPF problem.
\end{remark}
\subsection{Surrogate Model for Voltage Stability Index}
An alternative feasible approach is to establish a high-precision surrogate model to accurately express the VSI. Owing to its strong fitting ability, the Back Propagation NN (BPNN) is applied to obtain the approximate analytic expression of $\sigma$ in this study. 
Before proceeding with the establishment of the NN-based surrogate model, we first introduce two variable sets: The first set $\boldsymbol{y}\coloneqq\{p_{g,i},v_i\mid i\in\mathcal{N}_{PV,ref}\}$ contains all control variables and the other set $\boldsymbol{z}\coloneqq\{v_i,\theta_j\mid i\in\mathcal{N},j\in\mathcal{N}_{PV,PQ}\}$ includes all voltage magnitudes and angles.  

A high-quality NN model hinges on a large quantity of properly sampled training data. To this end, the Latin hypercube sampling (LHS) method blended with rejection sampling is employed to generate sufficient eligible samples of $\boldsymbol{y}$. Subsequently, standard power flow analysis is carried out to determine the corresponding $\boldsymbol{z}$ and $\sigma$ samples.
Since the Jacobian $\boldsymbol{J}$ is directly related to voltage magnitudes and angles, this study adopts the samples of $\boldsymbol{z}$ and $\sigma$ as the voltage stability dataset $D=\{(\boldsymbol{z},\sigma)\}$. Detailed steps are summarized as follows:

\textit{Step 1:} Use LHS to generate abundant samples $\{\boldsymbol{y}^{(i)}\}_{i=1}^{M}$ within the hypercube defined by technical limits of all control variables.

\textit{Step 2:} Conduct power flow analysis at each $\boldsymbol{y}^{(i)}$ and reject all samples falling outside the feasible space, resulting in the feasible samples $\{\boldsymbol{y}^{(i)},\boldsymbol{z}^{(i)},\sigma^{(i)}\}_{i=1}^{N_s}$.

\textit{Step 3:} Train a BPNN model with $\boldsymbol{z}$ being predictors and $\sigma$ being the response. The explicit expression for $\sigma$ is given by a combination of function composition and matrix multiplication:
\begin{equation}\label{eq:6}
        \sigma(\boldsymbol{z})\coloneqq f^{L}(\boldsymbol{W}^{L}\boldsymbol{f}^{L-1}(\boldsymbol{W}^{L-1}\cdots\boldsymbol{f}^{1}(\boldsymbol{W}^{1}\boldsymbol{z})\cdots)),
\end{equation}
where $L$ is the number of layers, $\boldsymbol{W}^{L}=(w_{jk}^l)$ are weights between layer $l-1$ and $l$, and $\boldsymbol{f}^l$ is the activation function at layer $l$.

Although the VSC-OPF problem can be directly solved by the interior point method (IPM) after substituting $\sigma$ with \eqref{eq:6}, the inclusion of uncertain renewable generations introduces a challenging aspect of managing chance constraints. Moreover, the IPM requires the hessian of \eqref{eq:6}, which will be prohibitively time-consuming and memory-intensive for large-scale systems. These two issues will be addressed in subsequent sections.
\section{Practical Uncertainty Propagation Approach}\label{approach}
The main focus of this section is to provide a novel approach known as APCE that effectively deals with chance constraints and transcends the limitations of existing PCE methods.   
\subsection{Motivation}
Let the variable set $\boldsymbol{x}\coloneqq\{p_{g,i},q_{g, i},v_j,i_l,\sigma\}$ comprise all uncertainty-affected variables. The implicit functional dependency between $\boldsymbol{x}$ and the uncertain renewable generations $\boldsymbol{\xi}$ at a specified operation point $\boldsymbol{y}$ can be defined in the form of:
\begin{equation}
    \boldsymbol{x}=\boldsymbol{h}(\boldsymbol{y},\boldsymbol{\xi}).
\end{equation}

It is mathematically complex or even impossible to analytically derive the probability distribution of $\boldsymbol{x}$. In practice, the truncated PCE technique can be adopted as a surrogate model to approximate $x_i$ as a finite sum of polynomials of $\boldsymbol{\xi}$ \cite{feinberg2018multivariate}:
\begin{equation}
    \hat{x}_i=\sum\limits_{j=0}^Ma_j\boldsymbol{\Psi}_j(\boldsymbol{\xi}),
\end{equation}
where $\{\boldsymbol{\Psi}_j(\boldsymbol{\xi})\}$ is the set of multivariate orthogonal polynomial basis and $\{a_j\}$ is the set of unknown coefficients. The crux of PCE is the identification of the basis for dependent random variables. Although Nataf transformation has been adopted to transform dependent variables into independent ones \cite{xu2021iterative}, it assumes a Gaussian copula as the underlying dependence structure, which might not align with the real case. 
\begin{figure}[!t]
\subfloat[]{\includegraphics[width=1.65in]{Gaussian_copula.pdf}}
\quad
\subfloat[]{\includegraphics[width=1.65in]{APCE.pdf}}
\caption{Comparison between MC simulations and: (a) PCE; (b) APCE}
\label{fig:1}
\vspace{-15pt}
\end{figure}
This claim is exemplified in Fig. \ref{fig:1}, which compares the line flow distributions calculated with the MC simulations, the Gaussian-copula-based PCE, and the APCE that will be covered in the succeeding section. It is obvious that the proposed APCE outperforms the other as it accurately captures the propagated line flow distribution. 
\subsection{Adaptive Polynomial Chaos Expansion}\label{APCE}
This part will present how to construct the orthonormal polynomial basis for dependent variables through the APCE method.
It is important not to conflate the APCE with the generalized PCE in \cite{xiu2002wiener} and the arbitrary PCE in \cite{oladyshkin2012data} since the latter two approaches are both incapable of handling nonlinear dependence.
The salient feature of APCE lies in its capability to inherently accommodate the interdependence between different renewable generations and arbitrary marginal probability distributions under very mild assumptions. As a result, the APCE is highly flexible and adaptive, boasting decent accuracy in uncertainty propagation, as verified in Fig. \ref{fig:1}. Moreover, the proposed APCE can be implemented as entirely data-driven, eliminating the need for estimating either the marginal distribution or the joint probability distribution.
\subsubsection{Gram Matrix}
To start with, we first develop a set of monomial basis that can be subsequently combined into orthogonal polynomials.
Recall that given an $N$-dimensional multi-index $\boldsymbol{j}\coloneqq[j_1,\dots,j_N]\in\mathbb{N}^N$, a monomial of the random vector $\boldsymbol{\xi}$ is defined as the following product type:
\begin{equation}
    \boldsymbol{\xi}^{\boldsymbol{j}}=\xi_1^{j_1}\dots\xi_N^{j_N},
\end{equation}
which has a total degree $\vert\boldsymbol{j}\vert=j_1+\dots+j_N$.

Considering a fixed $u\in\mathbb{N}$, let $\mathcal{P}_u^N$ denote the multi-index set consisting of all $N$-dimensional multi-indices $\boldsymbol{j}\in\mathbb{N}^N$ such that $\vert\boldsymbol{j}\vert=u$. In other words, $\mathcal{P}_u^N$ is the set of all multi-indices with a fixed norm of $u$.
Then, it is elementary to show that the cardinality of $\mathcal{P}_u^N$ is $\binom{N+u-1}{u}$.
Subsequently, the cardinality of the multi-index set with the norm being at most $m$ can be calculated by summing up the cardinalities of the sets $\mathcal{P}_u^N$ for all $u\in\mathbb{N}$: 
\begin{equation}\label{eq:10}
    L_{m}=\sum\limits_{u=0}^m\binom{N+u-1}{u}=\binom{N+m}{m}.
\end{equation}
Denote by $\boldsymbol{j}^{(1)},\dots,\boldsymbol{j}^{(L_{m})}$ the sorted elements of this multi-index set using the graded reverse lexicographical ordering, then the corresponding monomials of $\boldsymbol{\xi}$ are defined by $\boldsymbol{\xi}^{\boldsymbol{j}^{(1)}},\dots,\boldsymbol{\xi}^{\boldsymbol{j}^{(L_{m})}}$, which also constitute the basis of the space of homogeneous polynomials with a degree at most $m$.

Gathering the above basis yields an $L_{N,m}$-dimensional monomial vector with the $i$-th element being $\boldsymbol{\xi}^{\boldsymbol{j}^{(i)}}$, i.e., $\boldsymbol{P}_m(\boldsymbol{\xi})=[\boldsymbol{\xi}^{\boldsymbol{j}^{(1)}},\dots,\boldsymbol{\xi}^{\boldsymbol{j}^{(L_{m})}}]^{\mathsf{T}}$. Thereafter, we further define the Gram matrix for monomials as follows.
\begin{definition}[Gram matrix]\label{def:1}
Given a random vector $\boldsymbol{\xi}$ and the corresponding monomial vector $\boldsymbol{P}_m(\boldsymbol{\xi})$, the Gram matrix is a matrix $\boldsymbol{G}_{m}$ whose $(p,q)$-th entry is the inner product of the $p$-th and $q$-th element of $\boldsymbol{P}_m(\boldsymbol{\xi})$:
\begin{equation}
    G_{m,pq}=\left\langle \boldsymbol{\xi}^{\boldsymbol{j}^{(p)}}, \boldsymbol{\xi}^{\boldsymbol{j}^{(q)}} \right\rangle=\int \boldsymbol{\xi}^{\boldsymbol{j}^{(p)}+\boldsymbol{j}^{(q)}}f_{\boldsymbol{\xi}}(\boldsymbol{\xi}) \, d\boldsymbol{\xi},
\end{equation}
where $f_{\boldsymbol{\xi}}(\boldsymbol{\xi})$ denotes the joint probability density function.
\end{definition}
In general, $f_{\boldsymbol{\xi}}(\boldsymbol{\xi})$ proves to be unavailable or challenging to estimate. To address this challenge, a data-driven method can be leveraged as an alternative:
\begin{equation}\label{eq:12}
    G_{m,pq}\approx\frac{1}{N_s}\sum\limits_{l=1}^{N_s}(\xi_1^{(l)})^{j_1^{(p)}+j_1^{(q)}}\dots(\xi_N^{(l)})^{j_N^{(p)}+j_N^{(q)}},
\end{equation}
where $\boldsymbol{\xi}^{(l)}=[\xi_1^{(l)},\dots,\xi_N^{(l)}]^{\mathsf{T}}$ denotes the $l$-th sample.
\subsubsection{Orthogonal Basis}
It is evident that $\boldsymbol{G}_m$ is positive-definite, thereby permitting the whitening transformation to generate orthogonal polynomials. 
In this paper, the Cholesky whitening matrix $\boldsymbol{W}_m$ is adopted which satisfies the condition $\boldsymbol{W}_m^{\mathsf{T}}\boldsymbol{W}_m=\boldsymbol{G}_m^{-1}$. This approach entails computing the Cholesky decomposition of $\boldsymbol{G}_m$ to obtain a lower-triangular matrix $\boldsymbol{U}_m$, which satisfies $\boldsymbol{G}_m=\boldsymbol{U}_m\boldsymbol{U}_m^\mathsf{T}$. The inverse of $\boldsymbol{U}_m$ then serves as a whitening matrix $\boldsymbol{W}_m$, which yields an $L_{m}$-dimensional polynomial vector:
\begin{equation}\label{eq:13}
    \boldsymbol{\Psi}_m(\boldsymbol{\xi})=\boldsymbol{W}_m\boldsymbol{P}_m(\boldsymbol{\xi})=\boldsymbol{U}_m^{-1}\boldsymbol{P}_m(\boldsymbol{\xi}).
\end{equation}
Next, we would like to state the appealing property of this new set of polynomials, as presented in Proposition \ref{prop:1}.
\begin{prop}[Orthonormality]\label{prop:1}
For any distinct $i,j\in\mathbb{N}\cap[1,L_{m}]$, the $i$-th and $j$-th element of the polynomial vector $\boldsymbol{\Psi}_m(\boldsymbol{\xi})$ are mutually orthonormal w.r.t. the probability measure of $\boldsymbol{\xi}$.
\end{prop}
\begin{proof}
With a slight abuse of notation, let $\Psi_i(\boldsymbol{\xi})$ denote the $i$-th component of the vector $\boldsymbol{\Psi}_m(\boldsymbol{\xi})$. To prove the orthonormality between $\Psi_i(\boldsymbol{\xi})$ and $\Psi_j(\boldsymbol{\xi})$, it suffices to show $\left\langle \Psi_i(\boldsymbol{\xi}), \Psi_j(\boldsymbol{\xi}) \right\rangle_{\boldsymbol{\xi}}=\delta_{ij}$, where $\delta_{ij}$ is the Kronecker delta. 
From Definition \ref{def:1}, it can be seen that $\boldsymbol{G}_m=\mathbb{E}\left[\boldsymbol{P}_m(\boldsymbol{\xi})\boldsymbol{P}_m^{\mathsf{T}}(\boldsymbol{\xi})\right]$. Hence, the second-order moment of $\boldsymbol{\Psi}_m(\boldsymbol{\xi})$ can be expressed and simplified by
\begin{equation}\label{eq:14}
\begin{split}
    \mathbb{E}\left[\boldsymbol{\Psi}_m(\boldsymbol{\xi})\boldsymbol{\Psi}_m^{\mathsf{T}}(\boldsymbol{\xi})\right]
    & = \mathbb{E}\left[\boldsymbol{W}_m\boldsymbol{P}_m(\boldsymbol{\xi})\boldsymbol{P}_m^{\mathsf{T}}(\boldsymbol{\xi})\boldsymbol{W}_m^{\mathsf{T}}\right] \\
    & = \boldsymbol{W}_m\mathbb{E}\left[\boldsymbol{P}_m(\boldsymbol{\xi})\boldsymbol{P}_m^{\mathsf{T}}(\boldsymbol{\xi})\right]\boldsymbol{W}_m^{\mathsf{T}} \\
    & = \boldsymbol{W}_m\boldsymbol{G}_m\boldsymbol{W}_m^{\mathsf{T}} = \boldsymbol{I}_m,
    \end{split}
\end{equation}
where $\boldsymbol{I}_m$ is the $L_{m}$-order identity matrix. Therefore, the inner product between $\Psi_i(\boldsymbol{\xi})$ and $\Psi_j(\boldsymbol{\xi})$ is
\begin{equation}
\begin{split}
        \left\langle \Psi_i(\boldsymbol{\xi}), \Psi_j(\boldsymbol{\xi}) \right\rangle_{\boldsymbol{\xi}}& = \int \Psi_i(\boldsymbol{\xi})\Psi_j(\boldsymbol{\xi})f_{\boldsymbol{\xi}}(\boldsymbol{\xi}) \, d\boldsymbol{\xi}\\
        & = \mathbb{E}\left[\Psi_i(\boldsymbol{\xi})\Psi_j(\boldsymbol{\xi})\right]=\delta_{ij},
\end{split}
\end{equation}
which completes the proof.
\end{proof}\vspace{-8pt}
\begin{remark}
The whitening transformation is fundamentally distinct from the isoprobabilistic transformation, such as the Rosenblatt transformation. The latter approach aims to convert dependent random variables into independent ones, relying on full knowledge of the joint probability distribution. In contrast, the whitening transformation represents a fresh perspective that utilizes a linear combination of the monomial basis to generate an orthonormal polynomial basis w.r.t. the original dependent probability measure $f_{\boldsymbol{\xi}}(\boldsymbol{\xi})$.
\end{remark}
\subsubsection{APCE Approximation}
Note that the dimension of $\boldsymbol{\Psi}_m$ is consistent with the polynomial space of degree at most $m$, implying that this set of orthogonal polynomials functions as a basis. As $m$ approaches infinity, the vector $\boldsymbol{\Psi}_m$ encompasses an infinite collection of orthonormal bases and thus enables the Fourier polynomial expansion of the random vector $\boldsymbol{x}$, which is the APCE, as rigorously stated in Theorem \ref{the:1}.
\begin{theorem}\label{the:1}
%Let $\left(\Omega,\mathcal{F},\mathbb{P}\right)$ be the probability space and the random vector $\boldsymbol{\xi}$ be the mapping from $\Omega$ to $\mathbb{R}^N$. 
Assume that the probability density function of $\boldsymbol{\xi}$ is compactly supported or is exponentially integrable. Given a set of multivariate orthonormal polynomials $\{\Psi_i(\boldsymbol{\xi}),1\leq i<\infty\}$ w.r.t. $f_{\boldsymbol{\xi}}(\boldsymbol{\xi})d\boldsymbol{\xi}$, then each finite variance computation model represented as a map $x_i=h_i(\boldsymbol{y},\boldsymbol{\xi})$ can be expanded as a Fourier series composed of $\{\Psi_i(\boldsymbol{\xi})\}$:  
\begin{equation}\label{eq:16}
    x_i=h_i(\boldsymbol{y},\boldsymbol{\xi})\sim \sum\limits_{j=1}^{\infty}a_{i,j}\Psi_j(\boldsymbol{\xi}),
\end{equation}
%where $a_{i,j},i=1,\dots,N,j=1,\dots,\infty$ are the expansion coefficients and 
where the symbol $\sim$ means the APCE of $x_i$ converges in mean-square, both in probability and distribution.
The coefficients $a_{i,j},i=1,\dots,N,j=1,\dots,\infty$ are defined as
\begin{equation}
    a_{i,j}=\int h_i(\boldsymbol{y},\boldsymbol{\xi})\Psi_j(\boldsymbol{\xi})f_{\boldsymbol{\xi}}(\boldsymbol{\xi}) \, d\boldsymbol{\xi}.
\end{equation}
\end{theorem}
%\renewcommand\qedsymbol{}
\begin{proof}
    Detailed proof can be found in \cite{rahman2018polynomial}.
\end{proof}\vspace{-5pt}
In practical applications, \eqref{eq:16} must be truncated into a finite sum via the total-degree truncation scheme. In this work, an $m$-th order APCE approximation comprises all polynomials of total degree less than or equal to $m$:
\begin{equation}\label{eq:18}
    \hat{x}_i=\sum\limits_{j=1}^{L_{m}}a_{i,j}\Psi_j(\boldsymbol{\xi}).
\end{equation}
The expansion coefficients can be calculated using a non-intrusive method based on ordinary least-squares regression, which is similar to other PCE methods and hence omitted here. Based on this established expansion, the empirical distribution of $x_i$ can thus be accessed inexpensively by feeding the renewable generation dataset $\{\boldsymbol{\xi}^{(i)}\}_{i=1}^{N_s}$ into \eqref{eq:18}.
\subsection{Dimensionally Decomposed APCE}\label{DD-APCE}
It should be acknowledged that the APCE, like most existing PCE methods, will succumb to the curse of dimensionality, thereby limiting its application to high-dimensional cases. Motivated by the sparsity-of-effects principle, which suggests that a lower-variate approximation suffices to capture the uncertainty of the outputs, we further adopt a novel DD-APCE method \cite{lee2023high} to restructure the basis functions of the APCE in a dimensionwise manner.

The heart of the DD-APCE lies in limiting the interaction order among input variables. Following this intuition, for each $m\in \mathbb{N}$ and $0\leq s\leq N$, we first introduce the reduced multi-index set, denoted by
\begin{equation}\label{eq:19}
    \begin{split}
        \Pi_{s,m}^{N}\coloneqq\{&\boldsymbol{j}=(\boldsymbol{j}_l, \mathbf{0}_{-l})\in\mathbb{N}^N:\boldsymbol{j}_l\in\mathbb{N}^{\vert l\vert}, \\
        &\vert l\vert\leq\vert \boldsymbol{j}_l\vert\leq m, 0\leq\vert l\vert\leq s\},
    \end{split}
\end{equation}
where $l\subseteq\{1,\dots,N\}$ is a subset and $-l\coloneqq\{1,\dots,N\}\backslash l$ is the corresponding complementary subset. Here, $\boldsymbol{j}=(\boldsymbol{j}_l,\mathbf{0}_{-l})$ denotes an $N$-dimensional multi-index whose $i$-th component is $j_i$ if $i\in l$ and 0 otherwise. Hence, $\boldsymbol{j}$ has $\vert l\vert$ nonzero elements. Since $l$ has at most $s$ elements, $\Pi_{s,m}^N$ contains all multi-indices with at most $s$ nonzero elements and with the norm being at most $m$. The cardinality of $\Pi_{s,m}^N$ is given by
\begin{equation}\label{eq:20}
    L_{s,m}\coloneqq 1+\sum\limits_{k=1}^s\binom{N}{k}\binom{m}{k}.
\end{equation}

With the graded reverse lexicographical ordering, the elements of $\Pi_{s,m}^N$ are arranged as $\boldsymbol{j}^{(1)},\dots,\boldsymbol{j}^{(L_{s,m})}$ and the corresponding monomials of $\boldsymbol{\xi}$ are denoted by $\boldsymbol{\xi}^{\boldsymbol{j}^{(1)}},\dots,\boldsymbol{\xi}^{\boldsymbol{j}^{(L_{s,m})}}$. Employing the same procedures as described in Section \ref{APCE}, we can calculate the Gram matrix $\boldsymbol{G}_{s,m}$ and apply a whitening transformation to this new set of monomials to generate a set of measure-consistent orthonormal polynomials $\boldsymbol{\Psi}_{s,m}(\boldsymbol{\xi})$. This, in turn, leads to an $s$-variate, $m$-th order DD-APCE approximation of the form:
\begin{equation}
    \hat{x}_i=\sum\limits_{j=1}^{L_{s,m}}a_{i,j}\Psi_j(\boldsymbol{\xi}).
\end{equation}

Compared to the regular APCE, the DD-APCE, with $L_{s,m}\leq L_{m}$, offers superior efficiency while maintaining acceptable accuracy. In particular, this advantage is expected to be more pronounced in high-dimensional cases. To illustrate, the number of basis functions in a second-order APCE approximation scales quadratically as $(N+2)(N+1)/2$, whereas that of the uni-variate DD-APCE grows only linearly as $2N+1$. 
Therefore, the DD-APCE is capable of generating measure-consistent orthonormal polynomials in a stable manner while exhibiting exceptional efficacy in high dimensions. These attributes render the DD-APCE a promising candidate for practical uncertainty quantification in power systems.
\section{The Integrated Data-driven Scheme}\label{scheme}
In the previous section, we discussed the characterization of model responses under complexly correlated uncertainties. Our current goal is to amalgamate all the aforementioned techniques into a data-driven scheme tailored to tackle CC-VSC-OPF problems. However, prior to accomplishing this objective, it is imperative to address the computational burden associated with \eqref{eq:6}, thereby enhancing the feasibility and practicality of the proposed approach.
\subsection{PLS-NN Framework}
Recall that the IPM necessitates the Hessian of \eqref{eq:6} when solving an ordinary VSC-OPF problem. Nevertheless, when dealing with a large-scale NN, deriving the analytical expression for the Hessian matrix can be a laborious and complex task. Furthermore, even the numerical evaluation of the Hessian matrix becomes impractical when the dimension $2N-1$ exceeds a certain threshold.

One workaround to overcome this deficiency is to compute the gradient and Hessian of a reduced set of latent variables $\boldsymbol{\upsilon}$ and then retrieve the derivatives w.r.t. the original predictor variables $\boldsymbol{z}$. To facilitate this inverse transform, linear dimension reduction techniques are preferred, which render the following generic formula:
\begin{equation}\label{eq:21}
    \boldsymbol{\upsilon}=\boldsymbol{g}(\boldsymbol{z})=\boldsymbol{z}\boldsymbol{P},
\end{equation}
where the dimensions of $\boldsymbol{\upsilon}, \boldsymbol{z}, \boldsymbol{P}$ are $1\times M, 1\times (2N-1)$ and $(2N-1)\times M$, respectively. The NN surrogate model is then trained with $\boldsymbol{\upsilon}$ as input instead of $\boldsymbol{z}$, and $\sigma$ as output, yielding a more parsimonious model $\hat{\sigma}(\boldsymbol{\upsilon})$. During each iteration of the IPM procedure, the evaluation of $\sigma(\boldsymbol{z})$, its gradient, and Hessian for a given value of $\boldsymbol{z}^{(i)}$ can be derived as:
\begin{subequations}\label{eq:22}
    \begin{alignat}{1}
            \sigma(\boldsymbol{z}^{(i)}) &=\hat{\sigma}(\boldsymbol{\upsilon}^{(i)})=\hat{\sigma}(\boldsymbol{z}^{(i)}\boldsymbol{P}),\\
        \nabla \sigma\mid_{\boldsymbol{z}^{(i)}}&=\nabla (\hat{\sigma}\circ \boldsymbol{g})\mid_{\boldsymbol{z}^{(i)}} \notag\\
        &=\boldsymbol{P}\nabla \hat{\sigma}\mid_{\boldsymbol{\upsilon}^{(i)}},\\
        \nabla^2\sigma\mid_{\boldsymbol{z}^{(i)}}&=\nabla^2 (\hat{\sigma}\circ \boldsymbol{g})\mid_{\boldsymbol{z}^{(i)}} \notag\\
        &=\boldsymbol{P}\nabla^2 \hat{\sigma}\mid_{\boldsymbol{\upsilon}^{(i)}} \boldsymbol{P}^{\mathsf{T}}.
    \end{alignat}
    \vspace{-6pt}
\end{subequations}
\begin{figure}[!t]
\removelatexerror
\begin{algorithm}[H]
    \caption{The PLS-NN Framework}
    Generate abundant samples $\{\boldsymbol{y}^{(i)}\}_{i=1}^{M}$ using LHS\;
    Obtain feasible samples $\{\boldsymbol{y}^{(i)},\boldsymbol{z}^{(i)},\sigma^{(i)}\}_{i=1}^{N_s}$ through power flow analysis \;
    \eIf{\rm the value of $2N-1$ is reasonable}
    {Train the NN-based surrogate model $\sigma(\boldsymbol{z})$ as shown in \eqref{eq:6}\;}
    {Choose the dimension of latent variables $\boldsymbol{\upsilon}$ \;
    Perform PLS on voltage stability dataset $D$ and obtain the transformation matrix $\boldsymbol{P}$\;
    Construct latent variables through \eqref{eq:21} and create dataset $\{(\boldsymbol{\upsilon},\sigma)\}$\;
    Train the NN-based surrogate model, which yields the reduced model $\hat{\sigma}(\boldsymbol{\upsilon})$.}
        \end{algorithm}
        \vspace{-10pt}
\end{figure}

As indicated by \eqref{eq:22}, the quality of the retrieved values, $\sigma(\boldsymbol{z})$, $\nabla \sigma(\boldsymbol{z})$ and $\nabla^2 \sigma(\boldsymbol{z})$, is contingent upon the prediction accuracy of the reduced model $\hat{\sigma}(\boldsymbol{\upsilon})$. To this end, PLS is proposed in this study to construct new predictor variables $\boldsymbol{\upsilon}$ \cite{de1993simpls}. As a supervised technique, PLS maximizes the predictive power of $\boldsymbol{\upsilon}$ by accounting for the variability in the response. By capturing the variables that are highly informative regarding the response, PLS serves as a more effective alternative to unsupervised methods like principal component analysis and is expected to attain a parsimonious surrogate model $\hat{\sigma}(\boldsymbol{\upsilon})$ with satisfactory accuracy. 
A detailed procedure for implementing the proposed PLS-NN framework is presented in Algorithm 1. Within this framework, the computational overheads related to VSC-OPF problems are significantly reduced since it only requires a $M \times M$ Hessian as opposed to the original $(2N-1)\times (2N-1)$ one, where the $M \times M$ Hessian can be calculated inexpensively via several differentiation techniques, such as central difference differentiation. Meanwhile, this workaround not only enhances the scalability of the proposed approach but also serves to mitigate the limitations of the NN model training under high-dimensional inputs.
\begin{figure}[!t]
    \removelatexerror
     \begin{algorithm}[H]
             \KwIn{$\{\boldsymbol{\xi}^{(i)}\}_{i=1}^{N_s}$ - $N$-dimensional renewable generation dataset}
        \KwOut{$\boldsymbol{y}$ - Final solution of CC-VSC-OPF}
        \caption{The Iterative Data-Driven Scheme}
        Fix renewable generations at their expected values and then obtain the surrogate model of the VSI, $\sigma(\boldsymbol{z})\backslash\hat{\sigma}(\boldsymbol{\upsilon})$, according to Algorithm 1\;
        Build the deterministic VSC-OPF model\;
        Specify parameters of DD-APCE: total degree $m$ and maximum interaction order $s$\;
        Create $L_{s,m}$-dimensional monomial basis $\boldsymbol{P}_{s,m}(\boldsymbol{\xi })$\;
        Construct approximated $L_{s,m}\times L_{s,m}$ Gram matrix $\boldsymbol{G}_{s,m}$ via \eqref{eq:12}\;
        Generate orthonormal polynomial basis $\boldsymbol{\Psi}_{s,m}(\boldsymbol{\upsilon})$ through whitening transformation\;
               Set $k=0$, and initialize $\boldsymbol{x}_{\min}^{(0)}$ and $\boldsymbol{x}_{\max}^{(0)}$\; 
        Solve the deterministic VSC-OPF and obtain the initial operation point, $\boldsymbol{y}^{(0)}$ \;
        \While{$k\leq k_{\max}$}
        {
                                          Calculate coefficients of the $s$-variate, $m$-th order DD-APCE metamodel\;
                Approximate the distributions of $\boldsymbol{x}$ through stochastic testing\;
                Calculate margins $\Delta\boldsymbol{x}_{\min/\max}^{(k)}$ via \eqref{eq:26}\;
                \eIf{$\Delta\boldsymbol{x}_{\min}^{(k)}=\mathbf{0}$ \& $\Delta\boldsymbol{x}_{\max}^{(k)}=\mathbf{0}$}
                {
                \textbf{break}\;
                }{
                Update upper and lower bounds via \eqref{eq:27}\;
                Update iteration count $k\gets k+1$\;
                Solve the deterministic VSC-OPF to update the operation point $\boldsymbol{y}^{(k)}$
                }
                           
        }
     \end{algorithm}
     \vspace{-10pt}
\end{figure}
\subsection{The Iterative Scheme}
So far, we have developed the PLS-NN framework to solve large-scale VSC-OPF problems and have devised the DD-APCE technique to evaluate chance constraints at a fixed operation point. Our next step involves employing an iterative scheme inspired by \cite{roald2018chance} to merge these techniques, which enables the efficient solution of CC-VSC-OPF problems. The whole procedure is presented in Algorithm 2.

Prior to commencing the iteration process, the implicit VSI is substituted with the NN-based surrogate model. Subsequently, the deterministic VSC-OPF problem is formulated by assigning the uncertain renewable generations their anticipated values. At each iteration $k$, the deterministic VSC-OPF problem is solved utilizing the IPM, where the upper and lower bounds of $\boldsymbol{x}$ are set as $\boldsymbol{x}_{\min}^{(k)}$ and $\boldsymbol{x}_{\max}^{(k)}$, respectively. Once the operation point $\boldsymbol{y}^{(k)}$ is determined, the expansion coefficients in \eqref{eq:18} are calculated, enabling the estimation of empirical distributions of $\boldsymbol{x}$. 
Then, the tightening margins $\Delta\boldsymbol{x}_{\min}^{(k)}$ and $\Delta\boldsymbol{x}_{\max}^{(k)}$, which are introduced to ensure the prescribed violation probabilities, are calculated in a heuristic way:  
\begin{subequations}\label{eq:26}
    \begin{align}
        \Delta x_{i,\min}^{(k)}&=\left[x_{i,\min}^{(0)}-Q_{\epsilon_i}^{(k)}\left(x_i\right)\right]^+,\\
        \Delta x_{i,\max}^{(k)}&=\left[Q_{1-\epsilon_i}^{(k)}\left(x_i\right)-x_{i,\max}^{(0)}\right]^+,
    \end{align}
\end{subequations}
where $Q_{\alpha}^{(k)}(x_i)$ is the $\alpha$ quantile of the estimated distribution of $x_i$ at iteration $k$.
At the end of each iteration, the upper and lower bounds are updated via
\begin{subequations}\label{eq:27}
    \begin{align}
        \boldsymbol{x}_{\min}^{(k+1)}&=\boldsymbol{x}_{\min}^{(k)}+\Delta\boldsymbol{x}_{\min}^{(k)},\\
        \boldsymbol{x}_{\max}^{(k+1)}&=\boldsymbol{x}_{\max}^{(k)}-\Delta\boldsymbol{x}_{\max}^{(k)}.
    \end{align}
\end{subequations}

Following the above transformation, the CC-VSC-OPF problem is reduced to the traditional AC OPF with an explicit voltage stability constraint, which is amenable to existing optimization solvers.
Moreover, it can be seen that this algorithm is entirely data-driven and distribution-free, rendering it highly suitable for practical applications.
\section{Simulations and Results}\label{simulation}
In this section, various case studies are conducted on the IEEE 14-bus, IEEE 30-bus, European 89-bus, and Illinois 200-bus systems whose data can be obtained from MATPOWER \cite{zimmer2011matpower}. Solar power plants are added to the systems to introduce uncertainties in the VSC-OPF problem. The hourly power output of solar stations is simulated through the web platform, Renewables.ninja \cite{pfenninger2016long}. 
All the simulations are implemented on a desktop with 3.80-GHz AMD Ryzen 7 5800X processors and 16 GB RAM. The program environment is MATLAB 2021b. The VSC-OPF problem is implemented with the \textit{callback} function in MATPOWER and is resolved with Ipopt.
\vspace{-10pt}
\subsection{Small test cases}
This case study is performed on the IEEE 14-bus and IEEE 30-bus systems, with the corresponding voltage stability thresholds being 0.56 and 0.24. Three 60 MW solar power plants and five 30 MW solar power plants are added to the above two systems, respectively. A total of 11680 output samples for each power plant are collected by selecting corresponding stations in the Zhoushan islands in China from 2015 to 2018. 
\subsubsection{Accuracy of NN-based Surrogate Model}
We first investigate the suitability of employing a NN regression model as a surrogate for the VSI. PLS is unnecessary since the scale of both two cases is manageable. A total number of 10000 $\{\boldsymbol{z},\sigma\}$ samples are generated, and the approximate analytical expression $\sigma(\boldsymbol{z})$ is obtained by training the NN using these simulated samples. Furthermore, an additional set of 10000 samples is randomly generated for the test process. The training and test results are displayed in Fig. \ref{fig:2}, which includes error histograms and linear regression of targets relative to net outputs. It can be seen that the prediction generated by \eqref{eq:6} tracks the target accurately across the whole range. In the 14-bus system, the mean absolute errors for the training and testing sets are $3.16\times10^{-6}$ and $3.34\times10^{-6}$, respectively. Likewise, the corresponding errors are $1.09\times10^{-6}$ and $1.12\times10^{-6}$ in the 30-bus system. These observations indicate that the trained neural network boasts excellent accuracy and strong generalization capability. %Therefore, it can serve as a surrogate model for representing the voltage stability index. 
\begin{figure}[!t]
\centering
\subfloat[Error histogram]{\includegraphics[width=1.74in]{NN14_hist.pdf}}
\subfloat[Error histogram]{\includegraphics[width=1.74in]{NN30_hist.pdf}}\\[-1ex]
\subfloat[Regression fit]{\includegraphics[width=1.74in]{NN14_reg.pdf}}
\subfloat[Regression fit]{\includegraphics[width=1.74in]{NN30_reg.pdf}}
\caption{Accuracy of NN-based surrogate model: left column for IEEE 14-bus system, and right column for IEEE 30-bus system}
\label{fig:2}
\vspace{-20pt}
\end{figure}
\subsubsection{Accuracy of Uncertainty Quantification}
Now, we provide empirical validation of the accuracy of the APCE in terms of uncertainty quantification, utilizing MC simulations with 11680 samples as the benchmark. In this case study, a second-order APCE is adopted to represent system responses. Specifically, we select the VSI and the reactive power generation as the quantities of interest (QoIs) for illustration purposes. For each system, we calculate the simulated probability distributions of QoIs under the same operation condition, employing kernel smoothing methods to derive distribution curves. Visual comparisons between the APCE and MC are presented in Fig. \ref{fig:3}. Consistent with the findings in other relevant works \cite{xu2021iterative,muhlpfordt2019chance}, the APCE metamodel of order 2 is sufficient to capture the uncertainty of state variables in power systems. Furthermore, the results suggest that the second-order APCE can effectively model the VSI, a more intricate quantity. In conclusion, the data-driven APCE can represent random system responses accurately without any presumptions or inferences regarding individual distributions and underlying dependence structure.
\begin{figure}[!t]
\centering
\subfloat[]{\includegraphics[width=1.74in]{APCE14_s.pdf}}
\subfloat[]{\includegraphics[width=1.74in]{APCE14_q.pdf}}\\[-1ex]
\subfloat[]{\includegraphics[width=1.74in]{APCE30_s.pdf}}
\subfloat[]{\includegraphics[width=1.74in]{APCE30_q.pdf}}
\caption{Comparison between MC simulations and APCE metamodel: top row for IEEE 14-bus system, and bottom row for IEEE 30-bus system}
\label{fig:3}
\vspace{-10pt}
\end{figure}
\subsubsection{Satisfaction of Chance Constraints}
Next, the following two CC-VSC-OPF methods are considered to illustrate the effectiveness of the proposed approach: 
(\textbf{M1}) The proposed \textbf{Algorithm 2};
(\textbf{M2}) The identical iterative scheme except that the chance constraints are evaluated via MC simulations in each iteration. 
The performances of \textbf{M1} and \textbf{M2} are tested in terms of the total cost, computational efficiency, and the maximum empirical constraint violation probability $\epsilon_J$ under different settings. The results are summarized in Table \ref{tab:1}. 
As anticipated, decreasing predetermined violation probabilities will entail an increase in the cost to ensure a safer operation. It can be observed from Table \ref{tab:1} that our proposed method, \textbf{M1}, yields solutions that are of comparable quality to that of \textbf{M2} across different settings while exhibiting superior computational efficiency, rendering itself a computationally enhanced algorithm for CC-VSC-OPF problems.   

Based on local testing, the average CPU time required to solve the OPF problem of 14-bus and 30-bus systems is 0.0170s and 0.0357s, respectively. When it comes to the VSC-OPF problem, this time increases to 0.363s and 4.746s. These appreciable differences suggest scalability issues of larger-scale systems, as will be discussed in Section \ref{large}.
\begin{table}[t!]
%\setlength\belowcaptionskip{-2ex}
\centering
\caption{Performances Of M1 and M2 Under Different Settings in IEEE 14-Bus and 30-Bus Systems}
\label{tab:1}
\begin{threeparttable}
\renewcommand{\arraystretch}{1}
\begin{tabular}{cccccc} 
\toprule
\multicolumn{6}{c}{\textbf{Group 1:}\;$\boldsymbol{\epsilon=5\%}$\quad \textbf{14-bus system}}                     \\ 
\midrule
Method & Iterations & Time [s] & OPF Time [s]\tnote{*} & Cost [\$] & $\epsilon_J$ [\%] \\ 
\hline
\textbf{M1}     & 16          & 13.24 & 5.84 & 5378.741  & 5.09                             \\ 
%\cline{5-5}
\textbf{M2}     & 21         & 419.29 & 7.71 & 5378.720  & 5.00                             \\ 
\toprule
\multicolumn{6}{c}{\textbf{Group 2:}\;$\boldsymbol{\epsilon=8\%}$\quad \textbf{14-bus system}}                       \\ 
\midrule
Method & Iterations & Time [s] & OPF Time [s] & Cost [\$] & $\epsilon_J$ [\%]                       \\ 
\hline
\textbf{M1}     & 2          & 1.69 & 0.69 & 5370.110  & 8.00                             \\
\textbf{M2}     & 3          & 60.42 & 1.04 & 5370.084  & 8.00                             \\ 
\toprule
\multicolumn{6}{c}{\textbf{Group 3:}\;$\boldsymbol{\epsilon=5\%}$\quad \textbf{30-bus system}}                     \\ 
\midrule
Method & Iterations & Time [s] & OPF Time [s]\tnote{*} & Cost [\$] & $\epsilon_J$ [\%] \\ 
\hline
\textbf{M1}     & 6          & 32.60 & 28.69 & 356.596  & 5.02                             \\ 
%\cline{5-5}
\textbf{M2}     & 10         & 264.80 & 50.08 & 356.598  & 5.00                             \\ 
\toprule
\multicolumn{6}{c}{\textbf{Group 4:}\;$\boldsymbol{\epsilon=8\%}$\quad \textbf{30-bus system}}                       \\ 
\midrule
Method & Iterations & Time [s] & OPF Time [s] & Cost [\$] & $\epsilon_J$ [\%]                       \\ 
\hline
\textbf{M1}     & 6          & 32.24 & 28.26 & 356.395  & 8.03                             \\
\textbf{M2}     & 12          & 316.84 & 59.44 & 356.394  & 8.00                             \\ 
\midrule
\end{tabular}
\begin{tablenotes}
            \item[*] Sum of OPF solving time in all iterations.
        \end{tablenotes}
\end{threeparttable}
\vspace{-15pt}
\end{table}
\subsubsection{Improvement of Voltage Stability}
This subsection highlights the merits of the proposed CC-VSC-OPF model in optimizing voltage stability. A comparative analysis is conducted by adopting the traditional CC-OPF model without the voltage stability constraint. Both models are solved under the same iterative scheme, with a confidence level of $1-\epsilon=0.95$. The resulting probability distributions of $\sigma$ are illustrated using the split violin plots. As vividly shown in Fig. \ref{fig:4}, prior to the incorporation of the voltage stability constraint, there exists a certain probability that $\sigma$ fails to meet the requisite threshold, which indicates the system's vulnerability to voltage collapse under specific contingencies. 
With the voltage stability constraint, the stability threshold can be satisfied at the predetermined confidence level, emphasizing the significance of integrating the voltage stability constraint within the optimization framework to avoid potential instability issues.
\begin{figure}
\centering
\includegraphics[width=3.3in]{Stability_improve_small.pdf}
\caption{Probability distribution of the VSI with and without voltage stability constraint in small cases}
\label{fig:4}
\vspace{-15pt}
\end{figure}
\vspace{-10pt}
\subsection{Real-world test cases}\label{large}
\begin{table*}[t]
    \centering
\caption{Performances Of M1 and M2 With Different PLS Components in 89-Bus and 200-Bus Systems}
\label{tab:2}
\renewcommand{\arraystretch}{1}
\begin{tabular}{cccccccc} 
\toprule
\multicolumn{1}{c|}{}&\multicolumn{4}{c}{89-bus system} & \multicolumn{3}{c}{200-bus system}\\
\multicolumn{1}{c|}{} &\shortstack{\textbf{M1}\\10 Components} &\shortstack{\textbf{M1}\\20 Components} &\shortstack{\textbf{M1}\\w/o PLS} &\multicolumn{1}{c|}{\textbf{M2}} &\shortstack{\textbf{M1}\\25 Components}&\shortstack{\textbf{M1}\\49 Components}&\shortstack{\textbf{M2}\\49 Components} \\
\midrule
\multicolumn{1}{c|}{Iteration} & 8 &8 &8 & \multicolumn{1}{c|}{8} &9&8&8\\
\multicolumn{1}{c|}{Time [s]} & 10.53 &12.20 &895.74 & \multicolumn{1}{c|}{1293.37} &32.21&59.01&1124.15\\
\multicolumn{1}{c|}{OPF Time [s]}&2.49&3.20&887.01&\multicolumn{1}{c|}{880.67} &7.50&38.40&39.72\\
\multicolumn{1}{c|}{$\epsilon_J$ [\%]} &5.02&5.00&5.02&\multicolumn{1}{c|}{5.00}&5.00&5.00&5.00\\
\multicolumn{1}{c|}{RMSE of $p$ [p.u.]}&0.0028&0.0029&8.37$\times10^{-6}$&\multicolumn{1}{c|}{0}&1.80$\times10^{-4}$&6.70$\times10^{-7}$&0 \\
\multicolumn{1}{c|}{RMSE of $v$ [p.u.]}&2.10$\times10^{-5}$&1.97$\times10^{-5}$&2.76$\times10^{-6}$&\multicolumn{1}{c|}{0}&0.0013&1.70$\times10^{-7}$&0 \\
\midrule
\end{tabular}
\vspace{-15pt}
\end{table*}
This case study is conducted on the European 89-bus and Illinois 200-bus systems, with the corresponding voltage stability thresholds being 2.05 and 0.15. Seven 750 MW solar power plants and ten 75 MW solar power plants are added to the above systems, respectively. 
\subsubsection{Effectiveness of PLS-NN Framework}
In this section, we will demonstrate the potential of the PLS-NN framework to overcome scalability issues. The accuracy of the PLS-NN framework is first assessed. To determine the dimension of $\boldsymbol{\upsilon}$ for each system, PLS is initially performed with $2N-1$ components. The percentage variance explained in the response $\sigma$ is plotted against the number of PLS components. As illustrated in Fig. \ref{fig:5}, approximately 99.9\% of the variance in $\sigma$ is explained by the first 20 components in the 89-bus system, while nearly 99.6\% of the variance is accounted for by the first 49 components in the 200-bus system. Accordingly, the dimensions of the latent variables $\boldsymbol{\upsilon}$ are set to 20 and 49 for the two respective systems. 
The performance of resulting reduced surrogate models $\hat{\sigma}(\boldsymbol{\upsilon})$ is presented in Fig. \ref{fig:6}, which demonstrates that the PLS-NN framework is capable of constructing a more parsimonious model while maintaining an acceptable level of accuracy. 
\begin{figure}[!t]
\vspace{-10pt}
\centering
\subfloat[89-bus system]{\includegraphics[width=1.74in]{89_variance.pdf}}
\subfloat[200-bus system]{\includegraphics[width=1.74in]{200_variance.pdf}}
\caption{Percent variance explained in the response}
\label{fig:5}
\vspace{-15pt}
\end{figure}
\begin{figure}[!t]
\centering
\subfloat[89-bus system]{\includegraphics[width=1.74in]{NN89_hist.pdf}}
\subfloat[200-bus system]{\includegraphics[width=1.74in]{NN200_hist.pdf}}
\caption{Accuracy of reduced surrogate model}
\label{fig:6}
\vspace{-15pt}
\end{figure}

Next, we aim to elucidate the computational advantages of the proposed framework. Specifically, a comprehensive comparative analysis is conducted, focusing on the solution results under varying numbers of PLS components. To evaluate the performance of our approach, the root mean square error (RMSE) of operation points is additionally calculated, using results obtained with the \textbf{M2} method as a benchmark. Note that in the 200-bus system, the deterministic VSC-OPF problem will not be solved within 8h without PLS due to the substantial computational burden involved in evaluating the 399$\times$399 Hessian matrix. Consequently, the benchmark results are calculated with 49 PLS components. As clearly contrasted in Table \ref{tab:2}, the inclusion of PLS significantly enhances computational efficiency. In the 89-bus system, the average CPU time for solving the OPF problem exceeds a hundred seconds without PLS. Conversely, upon incorporating PLS, this time diminishes to a mere fraction of a second. This efficiency gain is achieved by evaluating much smaller Hessian matrices (10$\times$10 or 20$\times$20), as opposed to the larger 177$\times$177 matrix. A comparable improvement is observed in the 200-bus system, wherein the overall CPU time can be reduced to less than one minute.

Moreover, by examining the RMSE of operation points obtained with different PLS components, it is evident that the solutions achieved with PLS are all in close proximity to the benchmark results. Besides, the distributions of the VSI before and after considering the voltage stability constraint are depicted in Fig. \ref{fig:7}, which verifies the capability of the proposed model to improve voltage stability.
This demonstrates that the proposed algorithm, when combined with PLS, can effectively handle the CC-VSC-OPF problem in large-scale systems while consistently delivering satisfactory results. 
\begin{figure}
\centering
\includegraphics[width=3.3in]{Stability_improve_large.pdf}
\caption{Probability distribution of the VSI with and without voltage stability constraint in real-world cases}
\label{fig:7}
\vspace{-15pt}
\end{figure}
\subsubsection{Comparison of Different Settings of DD-APCE}
As mentioned in Section \ref{DD-APCE}, the efficacy of the proposed APCE diminishes as the dimensionality of uncertainty increases. In this section, we advocate the DD-APCE when faced with high-dimensional random variables. Four different parameter settings are considered in the 200-bus system: (\textbf{S1}) $s=1,m=2$; (\textbf{S2}) $s=2,m=2$; (\textbf{S3}) $s=1,m=3$; (\textbf{S4}) $s=3,m=3$. Fig. \ref{fig:8} illustrates the probability distribution of $\sigma$ obtained under each setting, demonstrating that all four settings accurately characterize uncertainty. This conclusion is further supported by Table \ref{tab:3}, where the RMSE is computed against the benchmark result obtained with the \textbf{M2} method. While higher interaction orders yield more accurate results, they also entail a larger number of basis functions and, consequently, longer evaluation time. In contrast, in this case, settings \textbf{S1} and \textbf{S3} can expedite the determination of expansion coefficients by restricting the interaction order while maintaining sufficient accuracy. 
\begin{figure}[!t]
\centering
\includegraphics[width=3.4in]{Order_Comparison.pdf}
\caption{Probability distribution of the voltage stability index under different settings of DD-APCE}
\label{fig:8}
\vspace{-15pt}
\end{figure}
\begin{table}[!t]
    \centering
        \caption{Comparison of Different Parameter Settings of DD-APCE in the 200-bus system}
    \label{tab:3}
    \renewcommand{\arraystretch}{1}
    \begin{threeparttable}
    \begin{tabular}{c|cccc}
    \toprule
    &\textbf{S1}&\textbf{S2}&\textbf{S3}&\textbf{S4} \\
    \midrule
    No. of Basis Functions&21&66&31&286\\
    Evaluation Time [s]\tnote{*}&1.23&2.58&1.55&32.14\\
    $\epsilon_J$ [\%]&5.00&5.00&5.00&5.00\\
    RMSE of $p$ [$10^{-5}$p.u.]&1.3&0.067&0.34&0.045\\
    RMSE of $v$ [$10^{-6}$p.u.]&2.7&0.17&6.0&0.067\\
    \midrule
    \end{tabular}
    \begin{tablenotes}
            \item[*] The time to determine APCE coefficients per iteration.
        \end{tablenotes}
\end{threeparttable}
\vspace{-15pt}
\end{table}
The difference in the number of basis functions and evaluation time will become more pronounced as the dimension increases. For instance, when $N=50$, \textbf{S2} involves 1326 basis functions, whereas \textbf{S1} requires only 101.
Therefore, DD-APCE is expected to be markedly more effective in dealing with high-dimensional cases, provided that a lower-variate truncation is adequate for uncertainty propagation.
%\vspace{-15pt}
\section{Conclusion}\label{conclusion}
In this paper, a CC-VSC-OPF model is proposed to enhance system stability in the presence of arbitrarily distributed and complexly correlated uncertainties. To address the inherent intractability of this model, a two-fold strategy is employed. Firstly, the implicit voltage stability constraint is explicitly embedded by harnessing the NN-based surrogate model. Secondly, to avoid intensive numerical simulations, chance constraints are tackled in a two-step procedure, in which the first step is to quantity uncertainty at a fixed operation point using the APCE, and the second step is to iteratively update the operation point to meet desired requirements. Moreover, a PLS-NN framework, along with a variant of APCE, namely DD-APCE, is introduced to enhance the versatility of the proposed solution scheme in large-scale systems and high-dimensional uncertainty scenarios. This integration significantly expands the applicability of the scheme with only a minor compromise in accuracy. Through comprehensive case studies, the proposed solution scheme proves to be computationally enhanced and accurate enough for practical applications.

\ifCLASSOPTIONcaptionsoff
  \newpage
\fi



\bibliographystyle{IEEEtran}

\bibliography{main}
\end{document}


